\textbf{Задача 6} Докажите, что $A\in\np$ тогда и только тогда, когда существует такая машина $M$, работающая на логарифмической памяти, что 
\[x\in A\iff \exists s\ M(x,s)=1\]

\begin{proof}

Обозначим последний язык как $\textbf{K}$, цель $\np=\textbf{K}$.

Покажем, что $\textbf{K}\subset\np$

Пусть $A\in\textbf{K}$.

У машины $M$ полином от $|x|$ конфигураций, поэтому можно перестроить её, чтобы она работала на полиномиальной от $|x|$ памяти и за полиномиальное от $|x|$ время. Обозначим получившуюся машину за $M'$, тогда

\[x\in A\iff \exists s\ M(x,s)=1\iff \exists s\ M'(x,s)=1\]

Из последней равносильности следует, что $A\in\np$

Покажем, что $\np\subset\textbf{K}$

Покажем в начале, что $\mathsf{3SAT}\in\textbf{K}$. Там в качестве сертификата можно передавать значения переменных, а $M$ будет проверять, что каждый дизъюнкт выполнен, это можно сделать на логарифмической памяти.

Пусть $A\in\np$ произвольный, тогда существует машина $V(x,y)$ работающая за не более чем $t(|x|)$ времени и такая, что
\[x\in A\iff\exists y\ V(x,y)=1\]

Аналогично Теореме Кука-Левина
\[x\in A \iff \exists y\ \varphi(x, y)  \iff   \varphi(x, \cdot) \in \mathsf{3SAT}\]

Где $\varphi(x, y)$ формула в $3$-КНФ проверяющая, что $y_1, \ldots, y_{t(|x|)}$ корректная цепочка конфигураций $V$ и среди этих конфигураций есть принимающая (то есть $V$ возвращает $1$ при таком исполнении). 

$\varphi(x, y)$ логарифмически вычислима (так как машина $V$ заранее известна), поэтому $A\in\textbf{K}$. В силу произвольности $A$, $\np\subset\textbf{K}$.

А значит $\np=\textbf{K}$.
\end{proof}